
    %%%%%%%%%%%%%%%%%%%%%%%%%%%%%%%%%%%%%%%%%%%%%%%%%%%%%%%%%%%%%%%%%%%%%%%%%%%%%%%%
    %%% Classe do documento
    \documentclass[12pt, addpoints]{exam}
    \UseRawInputEncoding

    %%%%%%%%%%%%%%%%%%%%%%%%%%%%%%%%%%%%%%%%%%%%%%%%%%%%%%%%%%%%%%%%%%%%%%%%%%%%%%%%
    % preambulo.tex
    %
    % Arquivo de configuração de packages para uso o LaTeX, conforme minhas
    % preferências e modelos pessoais.
    %
    % Para maiores informações, visite:
    %    https://github.com/abrantesasf/latex
    %
    % NÃO ALTERE SE NÃO SOUBER O QUE ESTÁ FAZENDO!
    %%%%%%%%%%%%%%%%%%%%%%%%%%%%%%%%%%%%%%%%%%%%%%%%%%%%%%%%%%%%%%%%%%%%%%%%%%%%%%%%


    %%%%%%%%%%%%%%%%%%%%%%%%%%%%%%%%%%%%%%%%%%%%%%%%%%%%%%%%%%%%%%%%%%%%%%%%%%%%%%%%
    %%% Estruturas de controle:
    \input{/app/static/utils/controles}

    %%%%%%%%%%%%%%%%%%%%%%%%%%%%%%%%%%%%%%%%%%%%%%%%%%%%%%%%%%%%%%%%%%%%%%%%%%%%%%%%
    %%% Compilação condicional em PDF:
    \input{/app/static/utils/pdf}

    %%%%%%%%%%%%%%%%%%%%%%%%%%%%%%%%%%%%%%%%%%%%%%%%%%%%%%%%%%%%%%%%%%%%%%%%%%%%%%%%
    %%% Configurações de layout da página:
    \input{/app/static/utils/layout}

    %%%%%%%%%%%%%%%%%%%%%%%%%%%%%%%%%%%%%%%%%%%%%%%%%%%%%%%%%%%%%%%%%%%%%%%%%%%%%%%%
    %%% Configurações para autores:
    \input{/app/static/utils/autores}

    %%%%%%%%%%%%%%%%%%%%%%%%%%%%%%%%%%%%%%%%%%%%%%%%%%%%%%%%%%%%%%%%%%%%%%%%%%%%%%%%
    %%% Cabeçalho e rodapé:
    %%% Atenção! É MUITO DIFÍCIL ajustar apropriadamente os running headers
    %%% e running footers da primeira vez. Você pode confiar no LaTeX e deixar que
    %%% ele faça o ajuste automaticamente ou pode quebrar a cabeça e
    %%% tentar melhorar algo que o LaTeX já faz muito bem, mesmo que não fique
    %%% exatamente do jeito que você quer. O que você prefere? Se quiser ajustar
    %%% manualmente, descomente a linha abaixo e faça as configurações no arquivo
    %%% "utils/headings.tex".
    %\input{/app/static/utils/headings}

    %%%%%%%%%%%%%%%%%%%%%%%%%%%%%%%%%%%%%%%%%%%%%%%%%%%%%%%%%%%%%%%%%%%%%%%%%%%%%%%%
    %%% Configuração para as fontes do documento:
    %%% Se você quiser usar fontes próprias ou do sistema, precisará ajustar as
    %%% configurações no arquivo "utils/fontes.tex".
    %%% ATENÇÃO: por padrão eu utilizo XeLaTeX com fontes proprietárias projetadas
    %%% por Matthew Butterick, adquiridas em https://mbtype.com, que não podem ser
    %%% distribuídas devido à licença de utilização. Se você usar XeLaTeX, você
    %%% DEVERÁ OBRIGATORIAMENTE ajustar as fontes no arquivo "utils/fontes.tex".
    \input{/app/static/utils/fontes}

    %%%%%%%%%%%%%%%%%%%%%%%%%%%%%%%%%%%%%%%%%%%%%%%%%%%%%%%%%%%%%%%%%%%%%%%%%%%%%%%%
    %%% Sumário:
    \input{/app/static/utils/toc}

    %%%%%%%%%%%%%%%%%%%%%%%%%%%%%%%%%%%%%%%%%%%%%%%%%%%%%%%%%%%%%%%%%%%%%%%%%%%%%%%%
    %%% Matemática:
    \input{/app/static/utils/matematica}

    %%%%%%%%%%%%%%%%%%%%%%%%%%%%%%%%%%%%%%%%%%%%%%%%%%%%%%%%%%%%%%%%%%%%%%%%%%%%%%%%
    %%% Notas de rodapé:
    \input{/app/static/utils/footnotes}

    %%%%%%%%%%%%%%%%%%%%%%%%%%%%%%%%%%%%%%%%%%%%%%%%%%%%%%%%%%%%%%%%%%%%%%%%%%%%%%%%
    %%% Referências cruzadas, links, citações:
    \input{/app/static/utils/crossrefs}

    %%%%%%%%%%%%%%%%%%%%%%%%%%%%%%%%%%%%%%%%%%%%%%%%%%%%%%%%%%%%%%%%%%%%%%%%%%%%%%%%
    %%% Configurações para os hiperlinks em PDF:
    \input{/app/static/utils/pdfconfig}

    %%%%%%%%%%%%%%%%%%%%%%%%%%%%%%%%%%%%%%%%%%%%%%%%%%%%%%%%%%%%%%%%%%%%%%%%%%%%%%%%
    %%% Glossário e índice remissivo:
    \input{/app/static/utils/glossindex}

    %%%%%%%%%%%%%%%%%%%%%%%%%%%%%%%%%%%%%%%%%%%%%%%%%%%%%%%%%%%%%%%%%%%%%%%%%%%%%%%%
    %%% Referências bibliográficas:
    \input{/app/static/utils/bibliografia}

    %%%%%%%%%%%%%%%%%%%%%%%%%%%%%%%%%%%%%%%%%%%%%%%%%%%%%%%%%%%%%%%%%%%%%%%%%%%%%%%%
    %%% Cores:
    \input{/app/static/utils/cores}

    %%%%%%%%%%%%%%%%%%%%%%%%%%%%%%%%%%%%%%%%%%%%%%%%%%%%%%%%%%%%%%%%%%%%%%%%%%%%%%%%
    %%% Computação:
    \input{/app/static/utils/computacao}

    %%%%%%%%%%%%%%%%%%%%%%%%%%%%%%%%%%%%%%%%%%%%%%%%%%%%%%%%%%%%%%%%%%%%%%%%%%%%%%%%
    %%% Gráficos:
    \input{/app/static/utils/graficos}

    %%%%%%%%%%%%%%%%%%%%%%%%%%%%%%%%%%%%%%%%%%%%%%%%%%%%%%%%%%%%%%%%%%%%%%%%%%%%%%%%
    %%% Ícones:
    \input{/app/static/utils/icones}

    %%%%%%%%%%%%%%%%%%%%%%%%%%%%%%%%%%%%%%%%%%%%%%%%%%%%%%%%%%%%%%%%%%%%%%%%%%%%%%%%
    %%% Blocos:
    \input{/app/static/utils/blocos}

    %%%%%%%%%%%%%%%%%%%%%%%%%%%%%%%%%%%%%%%%%%%%%%%%%%%%%%%%%%%%%%%%%%%%%%%%%%%%%%%%
    %%% Boxes coloridos:
    \input{/app/static/utils/colorbox}

    %%%%%%%%%%%%%%%%%%%%%%%%%%%%%%%%%%%%%%%%%%%%%%%%%%%%%%%%%%%%%%%%%%%%%%%%%%%%%%%%
    %%% Tabelas:
    \input{/app/static/utils/tabelas}

    %%%%%%%%%%%%%%%%%%%%%%%%%%%%%%%%%%%%%%%%%%%%%%%%%%%%%%%%%%%%%%%%%%%%%%%%%%%%%%%%
    %%% Ambientes de listas:
    \input{/app/static/utils/listas}

    %%%%%%%%%%%%%%%%%%%%%%%%%%%%%%%%%%%%%%%%%%%%%%%%%%%%%%%%%%%%%%%%%%%%%%%%%%%%%%%%
    %%% Ambientes floats e similares:
    \input{/app/static/utils/floats}

    %%%%%%%%%%%%%%%%%%%%%%%%%%%%%%%%%%%%%%%%%%%%%%%%%%%%%%%%%%%%%%%%%%%%%%%%%%%%%%%%
    %%% Ajustes para a classe EXAM:
    \input{/app/static/utils/exam}

    %%%%%%%%%%%%%%%%%%%%%%%%%%%%%%%%%%%%%%%%%%%%%%%%%%%%%%%%%%%%%%%%%%%%%%%%%%%%%%%%
    %%% Ajustes para textos:
    \input{/app/static/utils/texto}

    %%%%%%%%%%%%%%%%%%%%%%%%%%%%%%%%%%%%%%%%%%%%%%%%%%%%%%%%%%%%%%%%%%%%%%%%%%%%%%%%
    %%% Ambientes, aliases, comandos e customizações em geral para serem
    %%% utilizadas nos documentos. Defina aqui qualquer customização específica
    %%% para seus documentos!
    \input{/app/static/utils/customizacoes}

    %%%%%%%%%%%%%%%%%%%%%%%%%%%%%%%%%%%%%%%%%%%%%%%%%%%%%%%%%%%%%%%%%%%%%%%%%%%%%%%%
    %%% Ajuste do layout, espaçamento de linhas e etc.:
    \geometry{a4paper, portrait, left=2.0cm, right=2.0cm, top=2.0cm, bottom=2.0cm}
    \singlespacing

    %%%%%%%%%%%%%%%%%%%%%%%%%%%%%%%%%%%%%%%%%%%%%%%%%%%%%%%%%%%%%%%%%%%%%%%%%%%%%%%%
    %%% Configurações de cabeçalho e rodapé (não use fancyhdr pois ocorre conflito):
    \pagestyle{headandfoot}
    \runningheadrule
    \firstpageheader{}{}{}
    \runningheader{Sem disciplina}{}{Avaliação}
    \firstpagefooter{}{}{Página \thepage\ de \numpages}
    \runningfooter{}{}{Página \thepage\ de \numpages}
    \runningfootrule

    %%%%%%%%%%%%%%%%%%%%%%%%%%%%%%%%%%%%%%%%%%%%%%%%%%%%%%%%%%%%%%%%%%%%%%%%%%%%%%%%
    %%% Configurações para as propriedades do PDF:
    \hypersetup{
    hidelinks,           % Comente para web, descomente para impressão
    colorlinks=true,      % True para web, False para impressão
    pdftitle=teste: Avaliação,
    pdfauthor=bre,
    pdfsubject=Sem disciplina,
    pdfkeywords=['Sem tag'],
    pdfinfo={
        CreationDate={}, % Ex.: D:AAAAMMDDHH24MISS
        ModDate={}       % Ex.: D:AAAAMMDDHH24MISS
    }
    }
    \input{/app/static/utils/pdfconfig}

    %%%%%%%%%%%%%%%%%%%%%%%%%%%%%%%%%%%%%%%%%%%%%%%%%%%%%%%%%%%%%%%%%%%%%%%%%%%%%%%%
    %%% Começa o documento
    \begin{document}

    %%%%%%%%%%%%%%%%%%%%%%%%%%%%%%%%%%%%%%%%%%%%%%%%%%%%%%%%%%%%%%%%%%%%%%%%%%%%%%%%
    %%% Folha de instruções padronizadas da UVV para a prova
    \pdfbookmark[1]{Instruções}{instrucoes}
    \begin{figure}[H]
        \begin{center}
            \includegraphics[scale=0.3]{{/app/static/imgs/logo-uvv.png}}
        \end{center}	
    \end{figure}

    \vspace{-1cm}

    \begin{table}[H]
        \centering
        \begin{tabular}{|llll|}
            \hline
            \multicolumn{2}{|l|}{\textbf{Disciplina:} Sem disciplina} & 
            \multicolumn{1}{l|}{\multirow{3}{*}{\textbf{\makecell{Nota:\\ \,\\ \,}}}} & 
            \multirow{3}{*}{\textbf{\makecell{Coordenador:\\ \,\\ \,}}} \\ 
            \cline{1-2}
            \multicolumn{2}{|l|}{\textbf{Professor:} bre \hspace{4.72cm} } & 
            \multicolumn{1}{l|}{} & \\ 
            \cline{1-2}
            \multicolumn{2}{|l|}{\textbf{Aluno:}} & 
            \multicolumn{1}{l|}{} & \\ 
            \hline
            \multicolumn{1}{|l|}{\textbf{Turma:} \hspace{6.3cm} } & 
            \multicolumn{1}{l|}{\textbf{Semestre:}} & 
            \multicolumn{2}{l|}{\textbf{Valor:} 10 pontos} \\ 
            \hline
            \multicolumn{1}{|l|}{\textbf{Data:}} & 
            \multicolumn{3}{l|}{\textbf{Avaliação:}} \\ 
            \hline
        \end{tabular}
    \end{table}

    \begin{center}
        \fbox{\setlength\fboxsep{0.3cm} \fbox{\parbox{15.4cm}{
            \begin{center}
                \textbf{INSTRUÇÕES PARA A REALIZAÇÃO DESTA AVALIAÇÃO:}
            \end{center}
            \begin{itemize}[leftmargin=*]
                \item Esta é a primeira avaliação da disciplina Sem disciplina, 
                    e engloba o conteúdo referente Sem tag.
                \item Esta avaliação contém 10 questões objetivas, \textbf{todas obrigatórias}.
                \item \textbf{Leia atentamente cada questão!} As questões objetivas terão uma,
                    e apenas uma única, resposta a ser marcada.
                \item \textbf{Desligue o celular} antes de começar. A avaliação é
                    \textbf{individual} e \textbf{sem consulta}.
                \item As questões podem ser respondidas, na prova, com lápis ou caneta. Ao final
                    da prova \textbf{{VOCÊ DEVERÁ PREENCHER O GABARITO COM CANETA
                    DE COR PRETA}} \textbf{\underline{OU AZUL}} para que seja feita a correção
                    automatizada através da leitura do padrão de respostas marcado no
                    gabarito.
                \item \textbf{CUIDADO ao preencher o gabarito} pois eles são \textbf{INDIVIDUAIS
                    e NOMEADOS} (cada estudante tem um gabarito próprio específico, com um
                    código de identificação único). \textbf{Questões rasuradas são anuladas}
                    pelo software de leitura do gabarito.
                \item Ao final da prova, \textbf{DEVOLVA A PROVA E O GABARITO} para o professor.
                \item Siga todas as normas de \textbf{Integridade Acadêmica} da disciplina.
                    Alunos que forem flagrados com qualquer espécie de ``cola'' ou trocando
                    informações com outros alunos terão suas avaliações recolhidas, as notas
                    zeradas, e a situação será encaminhada para a coordenação para a aplicação
                    das penalidades previstas pela universidade.
                \item Com 90 minutos de avaliação você terá tempo de até 9,min por questão,
                    em média.
                \item Boa avaliação!
            \end{itemize}
            \vspace{2cm}
        }}}
    \end{center}
    \newpage

    %%%%%%%%%%%%%%%%%%%%%%%%%%%%%%%%%%%%%%%%%%%%%%%%%%%%%%%%%%%%%%%%%%%%%%%%%%%%%%%%
    %%% Inicia o bloco de questões:
    \begin{questions}

    \newpage
    \question
Dentre as definições a seguir, ligadas ao conceito de visões do modelo relacional, qual delas é \textbf{INCORRETA}?

\begin{choices}
\choice Uma visão relacional é uma relação virtual, derivada de relações base a partir da especificação de operações da álgebra relacional.
\choice Uma visão é útil por representar uma percepção particular do banco de dados, compartilhado por muitos aplicativos.
\choice Uma visão relacional é uma relação virtual que nunca é materializada.
\choice O gerenciamento de visões envolve a conversão da consulta do usuário sobre as visões para a consulta sobre as relações base.
\choice Programas aplicativos do banco de dados podem ser executados sobre visões de relações da base de dados.
\end{choices}

\question
Uma prova de vestibular foi elaborada com 25 questões de múltipla escolha com 5 alternativas. O número de candidatos presentes à prova foi 63127. Considere a afirmação: Pelo menos 2 candidatos responderam de modo idêntico as \( k \) primeiras questões da prova. Qual é o maior valor de \( k \) para o qual podemos garantir que a afirmação é verdadeira.

\begin{choices}
\choice 9
\choice 8
\choice 6
\choice 10
\choice 7
\end{choices}

\question
(Adaptado do ENADE 2005) Analise o relacionamento existente entre as tabelas
``clubes' e ``jogadores', nas alternativas abaixo, e marque a alternativa que
indica que todo jogador deve pertencer a um único clube, e todo clube poder ter
ou não jogadores.
\begin{choices}
\choice \includegraphics[width=0.5\linewidth]{/app/static/imgs/defaul.png} \vspace{0.5cm}
\choice \includegraphics[width=0.5\linewidth]{/app/static/imgs/defaul.png} \vspace{0.5cm}
\choice \includegraphics[width=0.5\linewidth]{/app/static/imgs/defaul.png} \vspace{0.5cm}
\choice \includegraphics[width=0.5\linewidth]{/app/static/imgs/defaul.png} \vspace{0.5cm}
\choice \includegraphics[width=0.5\linewidth]{/app/static/imgs/defaul.png} \vspace{0.5cm}
\end{choices}

\question
Todos os convidados presentes num jantar tomam chá ou café. Treze convidados bebem café, dez bebem chá e 4 bebem chá e café. Quantas pessoas tem nesse jantar.

\begin{choices}
\choice 27
\choice 19
\choice 15
\choice 23
\choice 10
\end{choices}

\question
Sejam os seguintes predicados de uma linguagem de primeira ordem:

\begin{itemize}
    \item \( N(x) \): \( x \) é número;
    \item \( P(x) \): \( x \) tem propriedade \( P \);
    \item \( x < y \): \( x \) é menor que \( y \).
\end{itemize}

E sejam os símbolos:
\begin{itemize}
    \item \( \forall \): quantificador universal;
    \item \( \Rightarrow \): operador se-então;
    \item \( \neg \): operador de negação.
\end{itemize}

Para a fórmula: \( \forall x (N(x) \Rightarrow \neg \forall y (N(y) \Rightarrow y < x)) \), qual alternativa abaixo \textbf{NÃO} constitui uma tradução possível?

\begin{choices}
\choice Para qualquer \( x \), se \( x \) é número, então não é verdade que todos os números são menores do que ele.
\choice Não há um número tal que todos os números são menores do que ele.
\choice Para todo número, não é verdade que qualquer número seja menor do que ele.
\choice Para todo número, existe um outro número que é maior do que ele.
\choice Não há um número menor do que outro número.
\end{choices}

\question
Considere as seguintes tabelas em uma base de dados relacional:\\
\textbf{Departamento (CodDepto, NomeDepto)}\\
\textbf{Empregado (CodEmp, NomeEmp, CodDepto)}\\
Deseja-se obter uma tabela na qual cada linha é a concatenação de uma linha da tabela Departamento com uma linha da tabela de Empregado. Caso um departamento não possua empregados, sua linha no resultado deve conter vazio (NULL) nos campos referentes ao empregado. A operação de álgebra relacional que deve ser aplicada para combinar estas duas tabelas é:
\begin{choices}
\choice Divisão
\choice Junção interna
\choice União
\choice Projeção
\choice Junção externa
\end{choices}

\question
A arquitetura ilustrada na figura abaixo é um exemplo típico de qual sistema?
(CU = unidade de controle; IS = \ingles{stream} de instruções; PU = unidade de
processamento; DS = \ingles{stream} de dados; LM = memória local)
\begin{figure}[H]
\begin{center}
\includegraphics[width=0.5\linewidth]{/app/static/imgs/defaul.png}
\end{center}
\end{figure}
\vspace{-1cm}
\begin{choices}
\choice Processadores multicore
\choice Acesso não uniforme à memória
\choice Processadores seqüenciais
\choice Clusters
\choice Multiprocessadores simétricos
\end{choices}

\question
Seja a árvore binária abaixo a representação de um espaço de estados para um problema \( p \), em que o estado inicial é \( a \), e \( i \) e \( f \) são estados finais.

\begin{figure}[H]
    \centering
    \includegraphics[width=0.5\linewidth]{/app/static/imgs/defaul.png}
    \caption{Questão 59}
    \label{fig:enter-label}
\end{figure}

Um algoritmo de busca em largura-primeiro forneceria a seguinte sequência de estados como primeira alternativa a um caminho-solução para o problema \( p \):
\begin{choices}
\choice \( a \ b \ d \ e \ f \)
\choice \( a \ b \ c \ d \ e \ f \)
\choice \( a \ b \ d \ h \ e \ i \)
\choice \( a \ c \ f \)
\choice \( a \ b \ e \ i \)
\end{choices}

\question Assinale a alternativa \textbf{incorreta}:
\begin{choices}
\choice O modelo relacional de dados foi criado por Edgar Frank Cood
\choice Modelo de dados é a mesma coisa que o diagrama relacional
\choice No modelo relacional, tudo o que o usuário percebe são tabelas
\choice Um SGBD não tem interface gráfica
\choice O MariaDB é um SGBD
\end{choices}

\question
Um banco de dados OLAP (Online Analytical Processing --- Processamento Analítico
Online) é um banco de dados organizado de forma a dar suporte ao Business
Intelligence (BI --- Inteligência de Negócio) que, de forma simplificada, é o
processo de analisar dados para obter informações que podem ser utilizadas na
tomada de decisões comerciais e estratégicas. Em relação aos bancos de dados
OLAP, marque a afirmação correta:
\begin{choices}
\choice São otimizados para consulta e relatórios de rotina de dados.
\choice São preparados para responder perguntas do tipo ``Quais alunos estão matriculados na turma de Bancos de Dados I neste semestre?'
\choice São otimizados para o processamento de grande volume de dados históricos e métricas de negócio para a tomada de decisões estratégicas, permitindo e facilitando a análise de dados para a geração de informações e conhecimentos para a tomada de decisões.
\choice São otimizados para o processamento de grande volume de transações diárias, o que permite a análise e a tomada de decisões em tempo real, favorecendo assim a administração empresarial.
\choice Geralmente os dados de origem do OLTP são de bancos de dados OLAP.
\end{choices}



    %%%%%%%%%%%%%%%%%%%%%%%%%%%%%%%%%%%%%%%%%%%%%%%%%%%%%%%%%%%%%%%%%%%%%%%%%%%%%%%%
    %%% Finaliza o bloco de questões:
    \end{questions}
    \end{document}
    